\begin{figure}[H]
\centering
\begin{tikzpicture}[x=1cm, y=0.8cm, >=Latex]
  \draw[->] (-0.2,1.7) -- (4.7,1.7) node[right] {$\mathbb R$};
  \node[left] at (-0.2,1.7) {open};
  \foreach \r/\y in {4.0/2.35,2.8/2.0,1.6/1.65,0.8/1.3} {
    \draw[line width=1pt] (0,\y) -- (\r,\y);
    \draw (0,\y) circle (1.4pt);
    \draw (\r,\y) circle (1.4pt);
  }
  \node[below] at (0,1.05) {$0$ omitted};
  \node[below] at (2.5,1.05) {$(0,1/n)$};

  \begin{scope}[yshift=-2.4cm]
    \draw[->] (-0.2,1.7) -- (6.2,1.7) node[right] {$\mathbb R$};
    \node[left] at (-0.2,1.7) {unbounded};
    \foreach \l/\y in {0.4/2.35,1.3/2.0,2.2/1.65,3.1/1.3} {
      \draw[line width=1pt] (\l,\y) -- (5.6,\y);
      \fill (\l,\y) circle (1.4pt);
      \draw[->,line width=1pt] (5.4,\y) -- (5.9,\y);
    }
    \node[below] at (3.0,1.05) {$[n,\infty)$ escapes right};
  \end{scope}
\end{tikzpicture}
\caption{Closedness and boundedness are separate hypotheses: open intervals can lose their endpoint, and unbounded closed intervals can drift away.}
\label{fig:nested-interval-failures}
\end{figure}
